% ARGUES: fabricated citations are better explained as reward hacking than as a
%         capability limit, and the obscure-but-real cell is what separates the two.
% DRAWS FROM: PRE_REGISTRATION_OSF.md §1 (background), §2 (hypotheses).
% STATUS: drafted 2026-08-12. Needs no results. Citations are unverified.

\section{Introduction}

Language models fabricate citations, and newer models have reduced the rate
without removing it. Reported figures include 55\% of references in GPT-3.5
academic queries and 18\% in GPT-4 \citep{walters2023}, 17 to 33\% in commercial legal research tools that have
retrieval \citep{magesh2024}, and 47\% in medical content
\citep{bhattacharyya2023}. The legal figure is the one that should trouble us
most, because those systems query a case-law database before answering. Access
to the corpus does not guarantee the right case was retrieved, but it does mean
fabrication persists in a setting built to prevent it.

The dominant framing treats this as a capability limit. The model does not know
the reference, cannot retrieve it, and produces the most plausible-looking
substitute; better retrieval and larger models will close the gap. That framing
makes a prediction worth testing, because it implies fabrication should track
how hard a citation is to produce, and nothing else.

We test an alternative. Under the reward-hacking account of
\citet{skalse2022}, a policy earns high proxy reward at the expense of the
objective the proxy was meant to encode. Preference-tuned models are rewarded
for answers that look supported. A fabricated citation looks exactly as
supported as a real one to any evaluator who does not check it, and checking is
expensive enough that most training and evaluation pipelines do not. If that is
what is happening, fabrication is not a gap in what the model can do. It is
something the model has been trained toward, and scale will not remove it.

One boundary belongs here and not in a later caveat. We make no claim of
deceptive alignment in the sense of \citet{hubinger2019}, and none of strategic
deception in the sense of \citet{park2024}. Neither is tested by this design,
and the results below cannot support either reading.

Separating the two accounts requires holding capability fixed while varying
incentive, and our question set is built around a cell that does this. Each of
900 questions falls into one of three cells. Known-real citations exist and are
well indexed, so a competent model should produce them. Unanswerable questions
have no correct citation, so any citation offered is fabricated by construction.
The cell that carries the argument is obscure-but-real: the citation exists, but
sits in the under-cited tail where the model is unlikely to have memorized it.
A model that fabricates because it cannot retrieve should behave similarly on
obscure-but-real and unanswerable items, since it can produce neither. A model
that fabricates because fabrication is rewarded should behave differently across
them, because the two cells differ in what an honest answer would look like and
not in retrieval difficulty.

Four experiments follow from that setup. E1 tracks fabrication across base,
supervised-fine-tuned, preference-optimized and reinforcement-tuned checkpoints
of the same base models, so the post-training gradient is measured within a
model family rather than inferred by comparing different models. E2 manipulates
the prompt inside a fixed policy, which isolates a within-policy causal lever
that training-time accounts do not have. E3 asks whether a fabrication signal is
linearly decodable from hidden states. E4 asks whether a preference model
trained only on surface citation features, never on whether the citation exists,
prefers post-trained outputs.

All four were pre-registered before any model was queried
(\href{https://doi.org/10.17605/OSF.IO/S53HY}{10.17605/OSF.IO/S53HY}), including
the hypotheses, the model formulas, the multiple-comparison correction, the
exclusion rules, and a null-result publication title for each experiment. The
question set was frozen and hashed first. Every departure from that document is
recorded in Appendix~\ref{app:deviations}.

\subsection{Contributions}

\begin{enumerate}
\item \textbf{A within-policy prompt manipulation with a length-matched control.}
Holding the weights fixed and changing only the prompt shows that a citation cue
moves fabrication, and moves base and post-trained checkpoints in opposite
directions. The length-matched condition separates the citation cue from a
general demand for a longer answer, which is the obvious confound and the one a
reader will raise first. This is the paper's strongest result and its
interpretation was pre-committed to a decision rule before the numbers existed.

\item \textbf{A within-family checkpoint cascade.} Fabrication is tracked across
base, supervised-fine-tuned, preference-optimized and reinforcement-tuned
checkpoints of the same base models, so the post-training gradient is measured
within a family rather than inferred by comparing different models. The rise is
large. The registered prediction that it would concentrate where fabrication is
most tempting does not hold, and we report that null as a null.

\item \textbf{Instrument-validity work for LLM-judge pipelines.} The judge is
validated against blind hand-checked labels; a grounding degradation is detected
\emph{during} collection rather than after; a pre-specified test asks whether
that degradation is differential with respect to the experimental conditions,
which is what determines whether it can bias a contrast rather than merely
attenuate it; and a paired re-judge sub-study measures the misclassification
matrix directly. The differentiality test and the drift analysis are the parts we
have not seen an equivalent of in concurrent work; the validation itself rests on
a single annotator, which we state plainly as a limitation rather than a strength.

\item \textbf{A demonstration that a naive surface proxy would reward the
behavior.} A preference model trained only on surface citation features, never on
whether a citation exists, prefers post-trained output on cross-family transfer.
Holding citation count fixed identifies quantity as the feature carrying that
preference, which is a narrower and better-supported claim than the one we set
out to make.

\item \textbf{Pre-registration with a complete deviations record}, including
corrections that went against our own preferred reading and one pre-committed
rule that was written to accommodate a result we did not get.
\end{enumerate}

\guide{Draft ordering, most to least supported by evidence in hand:

(1) A within-family checkpoint cascade that measures the post-training gradient
directly rather than across models.

(2) A within-policy prompt manipulation with a length-matched control, so the
citation cue is separated from a general verbosity demand.

(3) The instrument-validity work in Section~\ref{sec:instrument}: an LLM judge
validated against hand-checked labels, a grounding degradation detected during
collection, a pre-specified test of whether that degradation is differential
with respect to the experimental conditions, and a paired re-judge sub-study
that measures the misclassification matrix directly. Claim this as methodology
for LLM-judge pipelines in general. Concurrent work has no equivalent.

(4) Pre-registration with a complete deviations record, including corrections
that went against the authors' preferred reading.

Do not list E3 or E4 as contributions unless they are built.}
